\section{Discussion}\label{sec:discussion}

\paragraph{What the benchmarks measure}
Every corpus in our suite can be told apart from every other with high accuracy (proxy $\mathcal{A}$-distance \PADmin--\PADmax{} on a scale whose maximum is 2). On the article corpora a model of surface statistics recovers nearly all of the in-domain signal, and on average more than half of it. The article and tweet corpora reach most of their final AUC with 64 training examples, and the heaviest feature weights point to outlet signatures and to the time window in which the fake class was collected (Section~\ref{sec:shortcuts}). To a large extent, then, in-domain scores measure how well a model recognises the \emph{collection process} behind each class. \citet{verma2026audit} reached the same conclusion for ISOT alone. Our matrix shows that it holds across units and constructs, and that what is lost on a new corpus is mostly ranking, with calibration a smaller part.

\paragraph{What to do about it}
The decomposition separates a cheap problem from a hard one. The operating-point part is cheap. Rate transfer needs no target labels and one line of code, and it recovered \RecRank{} of that loss for \textsc{tfidf-lr} and between \RecRankMin{} and \RecRankMax{} across detectors, while the textbook label-shift correction made things worse for all of them. Someone deploying a detector on a new feed should therefore fix the \emph{share} of items they flag, not a score threshold. If the base rate of the feed is expected to differ from the training data, labelling a small sample to estimate it is the better investment, since neither label-free prior estimator we tested could do it (Section~\ref{sec:op}).

The ranking part is hard. Source diversity helped: pooled training reached \LocoPooled{} held-out AUC against \LocoMeanSingle{} for a single source. But label-free invariance methods made ranking worse, and the pooled model is still far below the in-domain ceiling of \LocoCeiling{}. Invariance to corpus identity is not enough when the label itself is tied to how each corpus was put together. Better ranking will need better data, in particular corpora whose real and fake items come from the same outlets and time windows, or evidence-based detectors that do not lean on style \citep{pelrine2023gpt4,thibault2025guide}. Multi-domain architectures \citep{nan2021mdfend,zhu2022endef} should be evaluated on held-out corpora with this entanglement in mind.

\paragraph{Style as construct and as shortcut}
The Bengali results add a distinction that the English matrix alone would not show. Clickbait detection held up on unseen publishers and partly across a change of language, while veracity detection did not survive a change of corpus. A plausible explanation is that a style-sensitive model transfers when style \emph{is} the construct, as with clickbait headlines, and fails when style is only correlated with the label inside one collection, as with fake news. With a single Bengali corpus and labels made in different ways, we offer this as a hypothesis for benchmark designers, not a settled result: they should say which kind of cue their labels license.

\begin{protocol}
For any claim that a misinformation detector generalises:
(1)~\emph{verify label direction} with signature probes, and treat a same-task AUC below 0.5 as a data error until shown otherwise;
(2)~\emph{deduplicate} within and across corpora with near-duplicate matching, and drop label-conflicting groups;
(3)~report at least one \emph{same-task cross-corpus} result, with a ranking metric (ROC-AUC or AP) next to F1;
(4)~report the \emph{decomposition}: F1 at the default, prior-matched and (cross-fitted) oracle thresholds;
(5)~run three cheap \emph{shortcut probes}: a surface-statistics model, a group-identity baseline where metadata exist, and a 64-example learning-curve point;
(6)~report a \emph{permutation control}, several seeds, and intervals that resample corpora when averaging over a matrix;
(7)~deploy with a \emph{rate-based} operating point unless target labels are available, and estimate the target base rate from a labelled sample if it is expected to change.
\end{protocol}

\subsection{Limitations and threats to validity}\label{sec:limits}
\emph{Models.} Six of our detectors are linear or tree-based on sparse, static-embedding and frozen RoBERTa features. The seventh fine-tunes only the upper four layers of RoBERTa-base, on 800 training items per corpus and three seeds, because the study ran on a two-core CPU; full fine-tuning on the complete splits might do better in-domain, though prior work found that it does not reliably help across corpora \citep{verma2026audit,delloglio2026comparison}. We did not prompt large language models, which earlier studies found more robust across corpora \citep{pelrine2023gpt4,delloglio2026comparison}. The decomposition applies unchanged to any detector that outputs a score, and testing it on those families is the most direct extension.
\emph{Corpora.} Labels differ in kind (claims, stories, outlets), which is why we lean on same-task pairs. FakeNewsNet is used through its public titles only. NELA-GT and BanFakeNews could not be obtained in our environment, so the Bengali analysis uses one Bengali corpus and a translate-test rather than a second native corpus. PHEME labels rumours, not falsehood, and serves only as an additional unit. The article corpora enter as cleaned working copies (Section~\ref{sec:data}). Deduplication thresholds (Jaccard $\geq$ 0.8) trade recall for precision, and looser thresholds would remove more items.
\emph{Statistics.} Twitter15/16 and PolitiFact test sets are small, so their cells have wide intervals (Table~\ref{tab:pairs}). We report pooled seed-bootstrap intervals per cell, resample corpora for matrix-level quantities, and avoid conclusions from single small cells. Seven corpora are still few, and the corpus-level intervals are correspondingly wide.
\emph{Operating points.} Rate transfer worked here partly because the base rates of our corpora are similar. The prior-shift experiment shows how it degrades when they are not; that experiment resamples existing test sets and does not reproduce every way a real feed can change.
\emph{Scope.} Our results concern text-only detection. Propagation-based detectors face related spurious-correlation problems on Twitter15/16 and PHEME \citep{wu2022spurious,mu2023time,pelrine2021simple}, which we did not re-evaluate.
\emph{Ethics.} We use only released text of public posts, with no user identifiers and no re-collection. Tweet content remains subject to the platform's terms, so we release code, per-cell results and saved scores but no post text.
